\chapter{Nonlocal transverse curvature}\label{ch:nonlocal-curvature}

\section{From the Wiener--Laguerre scalar to geometry}

Chapter~\ref{ch:wiener-laguerre-scalar} isolates the RH-equivalent scalar condition
\begin{equation}\label{eq:xf5-q-frontier}
Q_\Xi(x,y)>0,
\qquad
x\in\mathbb R,
\qquad
0<|y|<\frac12.
\end{equation}
The next step is an exact geometric reformulation of the same scalar.  Set
\[
F_x(y):=|\Xi(x+iy)|^2.
\]
Since $\Xi$ is real-entire,
\[
\overline{\Xi(x+iy)}=\Xi(x-iy),
\]
and direct differentiation yields
\begin{equation}\label{eq:xf5-curvature-identity}
\boxed{
\frac12 F_x''(y)
=
|\Xi'(x+iy)|^2
-\Ree\!\left(
\Xi(x+iy)\overline{\Xi''(x+iy)}
\right)
=Q_\Xi(x,y).
}
\end{equation}
Equation~\eqref{eq:xf5-curvature-identity} has status \status{EXACT}.  Therefore the XF-4 global sign frontier is equivalent to strict transverse convexity,
\begin{equation}\label{eq:xf5-convexity-frontier}
\boxed{
F_x''(y)>0
\quad
\text{for every }x\in\mathbb R,
\quad0<|y|<\frac12.
}
\end{equation}
The universal statement in Eq.~\eqref{eq:xf5-convexity-frontier} retains
\[
\status{OPEN\_RH\_EQUIVALENT\_CRITERION}.
\]

\section{Connection with the classical growth criterion}

The function $F_x$ is even in $y$, hence
\[
F_x'(0)=0.
\]
Strict convexity on $0<y<1/2$ therefore implies
\begin{equation}\label{eq:xf5-inner-growth}
F_x'(y)>0,
\qquad 0<y<\frac12.
\end{equation}
The outer region $y\ge 1/2$, corresponding to $\Ree s\ge1$, is controlled by the standard symmetric Hadamard/logarithmic-derivative argument.  Combining the inner and outer regions recovers the classical positivity/growth route associated with Hinkkanen and Lagarias \cite{Lagarias1999,Planat2026Theta}.

The solver records the same implication chain in four breakable stages:
\begin{enumerate}
\item \path{Q_Xi strict positivity};
\item equivalently, strict transverse convexity;
\item consequently, positive vertical growth;
\item with the outer-halfplane closure, the classical RH-equivalent growth criterion.
\end{enumerate}
The first globally quantified condition retains \status{OPEN}.  Thus XF-5 changes the representation of the proof obligation and preserves its evidential status.

\section{Differentiated theta-kernel representation}

The repository normalization is
\begin{equation}\label{eq:xf5-repo-fourier}
\Xi(z)=2\int_0^\infty\Phi(u)\cos(zu)\,\dd u.
\end{equation}
Define
\[
D(z)=\int_0^\infty\Phi(u)\cos(zu)\,\dd u=\frac12\Xi(z).
\]
In diagonal coordinates
\[
u=a+b,
\qquad
v=a-b,
\qquad
a>|b|,
\]
write
\[
M(a,b)=\Phi(a+b)\Phi(a-b).
\]
The symmetrized first-growth kernel can be represented as
\begin{equation}\label{eq:xf5-growth-kernel}
K_{x,y}(a,b)
=
\frac a2\cos(2xb)\sinh(2ya)
+
\frac b2\cos(2xa)\sinh(2yb).
\end{equation}
Differentiation gives the local curvature kernel
\begin{equation}\label{eq:xf5-curvature-kernel}
\boxed{
L_{x,y}(a,b)
:=\partial_y K_{x,y}(a,b)
=
a^2\cos(2xb)\cosh(2ya)
+
b^2\cos(2xa)\cosh(2yb).
}
\end{equation}
Under Eq.~\eqref{eq:xf5-repo-fourier}, the XF-4 scalar becomes
\begin{equation}\label{eq:xf5-nonlocal-q}
\boxed{
Q_\Xi(x,y)
=4\iint_{a>|b|}
M(a,b)L_{x,y}(a,b)\,\dd a\,\dd b.
}
\end{equation}
The differentiated-kernel identity and normalization bookkeeping have status \status{EXACT}.  The sign of the complete integral over the full admissible domain remains the active global obligation.

\section{Exact positive anchor at the central slice}

For $x=0$,
\begin{equation}\label{eq:xf5-xzero-kernel}
L_{0,y}(a,b)
=a^2\cosh(2ya)+b^2\cosh(2yb)>0
\end{equation}
throughout the interior domain $a>|b|$, $a>0$.  Together with the positive theta mass $M(a,b)$, this gives pointwise positivity of the XF-5 curvature integrand at the central slice.  It supplies a kernel-level version of the positive-sign anchor used in the Wiener argument of Chapter~\ref{ch:wiener-laguerre-scalar}.

\section{Nonlocality firewall}

Planat's 2026 theta-kernel analysis derives an exact longitudinal--transverse decomposition of the Xi growth derivative \cite{Planat2026Theta}.  The companion August 2026 preprint proves, for that phase-aligned decomposition, that the two sectors cancel to all algebraic orders and that universal independent block positivity fails: sufficiently far in the oscillatory regime at least one aligned block is negative \cite{Planat2026Nonlocal}.

The repository therefore records the following research firewall:
\begin{itemize}
\item Global correlated sign control --- \status{ACTIVE}.
\item Exact resummation preserving sector correlation --- \status{ACTIVE}.
\item Positive representation of the fully coupled integral --- \status{ACTIVE}.
\item Universal independent phase-aligned block positivity.\par\noindent Status: \status{LITERATURE\_NO\_GO}.
\item Finite block sampling --- \status{NUMERICAL\_DIAGNOSTIC}.
\end{itemize}
The no-go result is representation-specific.  It narrows the admissible strategy to globally correlated control of Eq.~\eqref{eq:xf5-nonlocal-q} or to a different exact representation that preserves the sign-carrying remainder.

\section{Executable audit surface}

The module \path{critical_axis.nonlocal_curvature} exposes
\begin{itemize}
\item \path{xi_transverse_curvature(x,y)}, implementing Eq.~\eqref{eq:xf5-curvature-identity};
\item \path{xi_transverse_curvature_direct(x,y)}, independently differentiating $|\Xi|^2$ for regression;
\item \path{theta_growth_kernel(x,y,a,b)};
\item \path{theta_curvature_kernel(x,y,a,b)}, implementing Eq.~\eqref{eq:xf5-curvature-kernel};
\item \path{theta_curvature_integrand}, evaluating one local contribution with fail-closed diagonal coordinates.
\end{itemize}

The corresponding tests lock the factor of two in Eq.~\eqref{eq:xf5-curvature-identity}, compare analytic and direct second derivatives, differentiate Eq.~\eqref{eq:xf5-growth-kernel} independently, verify the $x=0$ positive anchor, and reject invalid $a,b$ domains.

\section{XF-5 frontier}

The new proof surface is therefore the complete correlated curvature integral
\begin{equation}\label{eq:xf5-final-frontier}
\boxed{
\iint_{a>|b|}
M(a,b)L_{x,y}(a,b)\,\dd a\,\dd b>0
}
\end{equation}
for every real $x$ and every $0<|y|<1/2$.  Viable next steps are an analytic global lower bound, an exact correlation-preserving resummation, a positive-definite representation of the fully coupled integral, or a certified counterexample on the admissible domain.  The universal sign statement retains \status{OPEN\_RH\_EQUIVALENT\_CRITERION}.
